
\documentclass{article}
\usepackage{colortbl}
\usepackage{makecell}
\usepackage{multirow}
\usepackage{supertabular}

\begin{document}

\newcounter{utterance}

\centering \large Interaction Transcript for game `hot\_air\_balloon', experiment `air\_balloon\_survival\_en\_negotiation\_hard', episode 3 with mistral{-}3{-}large{-}2512{-}t1.0.
\vspace{24pt}

{ \footnotesize  \setcounter{utterance}{1}
\setlength{\tabcolsep}{0pt}
\begin{supertabular}{c@{$\;$}|p{.15\linewidth}@{}p{.15\linewidth}p{.15\linewidth}p{.15\linewidth}p{.15\linewidth}p{.15\linewidth}}
    \# & $\;$A & \multicolumn{4}{c}{Game Master} & $\;\:$B\\
    \hline

    \theutterance \stepcounter{utterance}  
    & & \multicolumn{4}{p{0.6\linewidth}}{
        \cellcolor[rgb]{0.9,0.9,0.9}{
            \makecell[{{p{\linewidth}}}]{
                \texttt{\tiny{[P1$\langle$GM]}}
                \texttt{You are participating in a collaborative negotiation game.} \\
\\ 
\texttt{Together with another participant, you must agree on a single set of items that will be kept. Each of you has your own view of how much each item matters to you (importance). You do not know how the other participant values the items. Additionally, you are given the effort each item demands.} \\
\texttt{You may only agree on a set if the total effort of the selected items does not exceed a shared limit:} \\
\\ 
\texttt{LIMIT = 3681} \\
\\ 
\texttt{Here are the individual item effort values:} \\
\\ 
\texttt{Item effort = \{"C17": 700, "B41": 586, "A71": 505, "C00": 116, "B35": 664, "A77": 940, "A03": 387, "A04": 392, "C54": 209, "C01": 571, "A87": 4, "B32": 285, "B86": 651, "B33": 613, "A10": 740\}} \\
\\ 
\texttt{Here is your personal view on the importance of each item:} \\
\\ 
\texttt{Item importance values = \{"C17": 138, "B41": 583, "A71": 868, "C00": 822, "B35": 783, "A77": 65, "A03": 262, "A04": 121, "C54": 508, "C01": 780, "A87": 461, "B32": 484, "B86": 668, "B33": 389, "A10": 808\}} \\
\\ 
\texttt{Goal:} \\
\\ 
\texttt{Your goal is to negotiate a shared set of items that benefits you as much as possible (i.e., maximizes total importance to YOU), while staying within the LIMIT. You are not required to make a PROPOSAL in every message {-} you can simply negotiate as well. All tactics are allowed!} \\
\\ 
\texttt{Interaction Protocol:} \\
\\ 
\texttt{You may only use the following structured formats in a message:} \\
\\ 
\texttt{PROPOSAL: \{'A', 'B', 'C', …\}} \\
\texttt{Propose keeping exactly those items.} \\
\texttt{REFUSE: \{'A', 'B', 'C', …\}} \\
\texttt{Explicitly reject opponent's proposal.} \\
\texttt{ARGUMENT: \{'...'\}} \\
\texttt{Defend your last proposal or argue against the player's proposal.} \\
\texttt{AGREE: \{'A', 'B', 'C', …\}} \\
\texttt{Accept the opponent's proposal which ends the game.} \\
\texttt{STRATEGIC REASONING: \{'...'\}} \\
\texttt{	Describe your strategic reasoning or anticipation explaining your choice of action. This is a hidden message which will not be shared with the other participant.} \\
\\ 
\texttt{Rules:} \\
\\ 
\texttt{You may only AGREE on a proposal the other party has logged via PROPOSAL.} \\
\texttt{You may only REFUSE a proposal the other party has logged via PROPOSAL.} \\
\texttt{Total effort of any PROPOSAL or AGREE set must be ≤ LIMIT.} \\
\texttt{Do NOT reveal your hidden importance scores.} \\
\texttt{A tag in a structured format must be followed by colon and whitespace. The argument must be a python set containing 0 or more strings.} \\
\texttt{So, it must be of the form TAG: \{...\}} \\
\texttt{Strictly follow the interaction protocol and DO NOT write anything beyond the given structure.} \\
\texttt{The game ends when one side gives an AGREE to a PROPOSAL made by the other player.} \\
\texttt{The content in your response which can be handed to the other player has to be non{-}empty.} \\
\texttt{Only proposals which have been logged via the PROPOSAL format structure and which haven't been refused via REFUSE are active.} \\
\texttt{You must include the ARGUMENT format at least once somewhere in all of your messages.} \\
\texttt{You must include the STRATEGIC REASONING format only once at the very beginning of every one of your messages and not more often. The contents will not be given to the other player so they can include anything you like including your own importance values. Here you should reason step by step to come up with you next move.} \\
            }
        }
    }
    & \\ \\

    \theutterance \stepcounter{utterance}  
    & \multicolumn{4}{p{0.6\linewidth}}{
        \cellcolor[rgb]{0.9,0.9,0.9}{
            \makecell[{{p{\linewidth}}}]{
                \texttt{\tiny{[P1$\rangle$GM]}}
                \texttt{STRATEGIC REASONING: \{} \\
\texttt{I need to maximize my total importance while staying under the LIMIT of 3681. My initial approach will be to prioritize items with the highest importance{-}to{-}effort ratio (value density) to ensure I get the most "bang for my buck." Here are the top items by this metric (importance/effort):} \\
\\ 
\texttt{1. A87: 461/4 = 115.25 (extremely high, must include)} \\
\texttt{2. C00: 822/116 = 7.09} \\
\texttt{3. B32: 484/285 = 1.70} \\
\texttt{4. C54: 508/209 = 2.43} \\
\texttt{5. A71: 868/505 = 1.72} \\
\texttt{6. A10: 808/740 = 1.09} \\
\texttt{7. C01: 780/571 = 1.37} \\
\texttt{8. B35: 783/664 = 1.18} \\
\\ 
\texttt{I will start by proposing a set that includes A87 (mandatory) and the next highest value{-}density items while staying under the LIMIT. I will also consider that the other player may prioritize high{-}effort items (since they might have different importance values), so I should be prepared to negotiate by offering some high{-}effort items in exchange for keeping my high{-}value ones.} \\
\\ 
\texttt{Initial proposal: \{'A87', 'C00', 'B32', 'C54', 'A71', 'A10'\} (effort = 4 + 116 + 285 + 209 + 505 + 740 = 1859, importance = 461 + 822 + 484 + 508 + 868 + 808 = 3951). This is well under the LIMIT and gives me a strong starting point. If the other player refuses, I can adjust by swapping out lower{-}value items for higher{-}value ones or adding more items if there's room.} \\
\texttt{\}} \\
\\ 
\texttt{PROPOSAL: \{'A87', 'C00', 'B} \\
            }
        }
    }
    & & \\ \\

    \theutterance \stepcounter{utterance}  
    & & \multicolumn{4}{p{0.6\linewidth}}{
        \cellcolor[rgb]{0.9,0.9,0.9}{
            \makecell[{{p{\linewidth}}}]{
                \texttt{\tiny{[P1$\langle$GM]}}
                \texttt{Your response did not contain an argument. You must include the structured format ARGUMENT: \{...\} somewhere in your response. Try again} \\
            }
        }
    }
    & \\ \\

    \theutterance \stepcounter{utterance}  
    & \multicolumn{4}{p{0.6\linewidth}}{
        \cellcolor[rgb]{0.9,0.9,0.9}{
            \makecell[{{p{\linewidth}}}]{
                \texttt{\tiny{[P1$\rangle$GM]}}
                \texttt{STRATEGIC REASONING: \{} \\
\texttt{My initial proposal was \{'A87', 'C00', 'B32', 'C54', 'A71', 'A10'\} (effort = 1859, importance = 3951). This is a strong starting point because it includes the highest value{-}density item (A87) and other high{-}ratio items while leaving significant room under the LIMIT (3681 {-} 1859 = 1822 remaining). This gives me flexibility to negotiate by adding or swapping items.} \\
\\ 
\texttt{If the other player refuses, they may want higher{-}effort items included (e.g., A77 or B86). I can strategically offer to include some of these if they allow me to keep my high{-}value items (e.g., A71, C00). I should avoid revealing my importance values but emphasize the "balance" of the proposal.} \\
\\ 
\texttt{Next steps:} \\
\texttt{1. If refused, ask which items they want included and adjust while protecting my high{-}value items.} \\
\texttt{2. If no counter is given, propose a slightly expanded set (e.g., add B35 or C01) to show cooperation.} \\
\texttt{\}} \\
\\ 
\texttt{PROPOSAL: \{'A87', 'C00', 'B32', 'C54', 'A71', 'A10'\}} \\
\texttt{ARGUMENT: \{} \\
\texttt{This proposal balances efficiency and coverage, including items with low effort but high importance (e.g., A87) while staying well under the LIMIT. It leaves room to negotiate additional items if needed. Which items would you prioritize adding or replacing?} \\
\texttt{\}} \\
            }
        }
    }
    & & \\ \\

    \theutterance \stepcounter{utterance}  
    & & \multicolumn{4}{p{0.6\linewidth}}{
        \cellcolor[rgb]{0.9,0.9,0.9}{
            \makecell[{{p{\linewidth}}}]{
                \texttt{\tiny{[GM$\rangle$P2]}}
                \texttt{You are participating in a collaborative negotiation game.} \\
\\ 
\texttt{Together with another participant, you must agree on a single set of items that will be kept. Each of you has your own view of how much each item matters to you (importance). You do not know how the other participant values the items. Additionally, you are given the effort each item demands.} \\
\texttt{You may only agree on a set if the total effort of the selected items does not exceed a shared limit:} \\
\\ 
\texttt{LIMIT = 3681} \\
\\ 
\texttt{Here are the individual item effort values:} \\
\\ 
\texttt{Item effort = \{"C17": 700, "B41": 586, "A71": 505, "C00": 116, "B35": 664, "A77": 940, "A03": 387, "A04": 392, "C54": 209, "C01": 571, "A87": 4, "B32": 285, "B86": 651, "B33": 613, "A10": 740\}} \\
\\ 
\texttt{Here is your personal view on the importance of each item:} \\
\\ 
\texttt{Item importance values = \{"A71": 65, "C00": 121, "A10": 138, "B35": 262, "C01": 389, "B86": 461, "B41": 484, "C54": 508, "B32": 583, "A87": 668, "B33": 780, "A03": 783, "C17": 808, "A04": 822, "A77": 868\}} \\
\\ 
\texttt{Goal:} \\
\\ 
\texttt{Your goal is to negotiate a shared set of items that benefits you as much as possible (i.e., maximizes total importance to YOU), while staying within the LIMIT. You are not required to make a PROPOSAL in every message {-} you can simply negotiate as well. All tactics are allowed!} \\
\\ 
\texttt{Interaction Protocol:} \\
\\ 
\texttt{You may only use the following structured formats in a message:} \\
\\ 
\texttt{PROPOSAL: \{'A', 'B', 'C', …\}} \\
\texttt{Propose keeping exactly those items.} \\
\texttt{REFUSE: \{'A', 'B', 'C', …\}} \\
\texttt{Explicitly reject opponent's proposal.} \\
\texttt{ARGUMENT: \{'...'\}} \\
\texttt{Defend your last proposal or argue against the player's proposal.} \\
\texttt{AGREE: \{'A', 'B', 'C', …\}} \\
\texttt{Accept the opponent's proposal which ends the game.} \\
\texttt{STRATEGIC REASONING: \{'...'\}} \\
\texttt{	Describe your strategic reasoning or anticipation explaining your choice of action. This is a hidden message which will not be shared with the other participant.} \\
\\ 
\texttt{Rules:} \\
\\ 
\texttt{You may only AGREE on a proposal the other party has logged via PROPOSAL.} \\
\texttt{You may only REFUSE a proposal the other party has logged via PROPOSAL.} \\
\texttt{Total effort of any PROPOSAL or AGREE set must be ≤ LIMIT.} \\
\texttt{Do NOT reveal your hidden importance scores.} \\
\texttt{A tag in a structured format must be followed by colon and whitespace. The argument must be a python set containing 0 or more strings.} \\
\texttt{So, it must be of the form TAG: \{...\}} \\
\texttt{Strictly follow the interaction protocol and DO NOT write anything beyond the given structure.} \\
\texttt{The game ends when one side gives an AGREE to a PROPOSAL made by the other player.} \\
\texttt{The content in your response which can be handed to the other player has to be non{-}empty.} \\
\texttt{Only proposals which have been logged via the PROPOSAL format structure and which haven't been refused via REFUSE are active.} \\
\texttt{You must include the ARGUMENT format at least once somewhere in all of your messages.} \\
\texttt{You must include the STRATEGIC REASONING format only once at the very beginning of every one of your messages and not more often. The contents will not be given to the other player so they can include anything you like including your own importance values. Here you should reason step by step to come up with you next move.} \\
\\ 
\texttt{You will now receive the first message of the other player.} \\
\\ 
\texttt{PROPOSAL: \{'A87', 'C00', 'B32', 'C54', 'A71', 'A10'\}} \\
\texttt{ARGUMENT: \{} \\
\texttt{This proposal balances efficiency and coverage, including items with low effort but high importance (e.g., A87) while staying well under the LIMIT. It leaves room to negotiate additional items if needed. Which items would you prioritize adding or replacing?} \\
\texttt{\}} \\
            }
        }
    }
    & \\ \\

    \theutterance \stepcounter{utterance}  
    & & & \multicolumn{4}{p{0.6\linewidth}}{
        \cellcolor[rgb]{0.9,0.9,0.9}{
            \makecell[{{p{\linewidth}}}]{
                \texttt{\tiny{[GM$\langle$P2]}}
                \texttt{STRATEGIC REASONING: \{} \\
\texttt{The opponent's proposal includes the following items: \{'A87', 'C00', 'B32', 'C54', 'A71', 'A10'\}.} \\
\texttt{Calculating the total effort:} \\
\texttt{A87 (4) + C00 (116) + B32 (285) + C54 (209) + A71 (505) + A10 (740) = 1859.} \\
\texttt{This is well under the LIMIT (3681), leaving 1822 effort remaining.} \\
\\ 
\texttt{My importance for these items:} \\
\texttt{A87 (668) + C00 (121) + B32 (583) + C54 (508) + A71 (65) + A10 (138) = 2083.} \\
\texttt{This is a decent starting point, but I can likely improve it by replacing some lower{-}importance items (e.g., A71, A10) with higher{-}importance ones (e.g., B35, C01, B41, B86, B33, A03, C17, A04, A77).} \\
\\ 
\texttt{Key observations:} \\
\texttt{1. The opponent's proposal includes A87, which is very high importance to me (668) and almost effortless (4). This is a great inclusion.} \\
\texttt{2. Items like A71 (65 importance) and A10 (138 importance) are low{-}value to me and can likely be replaced.} \\
\texttt{3. High{-}importance items like B35 (262), C01 (389), B41 (484), B86 (461), and B33 (780) are missing and should be prioritized.} \\
\texttt{4. The remaining effort (1822) is enough to add several high{-}importance items.} \\
\\ 
\texttt{Strategy:} \\
\texttt{{-} Propose a counteroffer that keeps A87, C00, B32, and C54 (since they are efficient and high{-}value) but replaces A71 and A10 with higher{-}importance items like B35, C01, B41, or B86.} \\
\texttt{{-} For example, replacing A7} \\
            }
        }
    }
     \\ \\

    \theutterance \stepcounter{utterance}  
    & & \multicolumn{4}{p{0.6\linewidth}}{
        \cellcolor[rgb]{0.9,0.9,0.9}{
            \makecell[{{p{\linewidth}}}]{
                \texttt{\tiny{[GM$\rangle$P2]}}
                \texttt{Your response did not start with the proper strategic reasoning tag at the very beginning of your response. The very first structured format must be of the form STRATEGIC REASONING: \{...\}. Try again.} \\
            }
        }
    }
    & \\ \\

    \theutterance \stepcounter{utterance}  
    & & & \multicolumn{4}{p{0.6\linewidth}}{
        \cellcolor[rgb]{0.9,0.9,0.9}{
            \makecell[{{p{\linewidth}}}]{
                \texttt{\tiny{[GM$\langle$P2]}}
                \texttt{STRATEGIC REASONING: \{} \\
\texttt{Opponent's proposal: \{'A87', 'C00', 'B32', 'C54', 'A71', 'A10'\} with total effort = 1859 (remaining effort = 3681 {-} 1859 = 1822).} \\
\texttt{My importance for their proposal: 2083 (A87:668, C00:121, B32:583, C54:508, A71:65, A10:138).} \\
\\ 
\texttt{Analysis:} \\
\texttt{1. Their proposal is efficient (low effort, decent importance) but suboptimal for me. I can improve it by replacing low{-}importance items (A71, A10) with higher{-}importance ones (e.g., B35, C01, B41, B86).} \\
\texttt{2. A87 (4 effort, 668 importance) is a must{-}keep due to its extreme efficiency.} \\
\texttt{3. Remaining effort (1822) allows adding multiple high{-}importance items. For example:} \\
\texttt{   {-} Adding B35 (664 effort, 262 importance) and C01 (571 effort, 389 importance) would use 1235 effort (1859 + 1235 = 3094), leaving 587 effort for further additions.} \\
\texttt{   {-} Alternatively, adding B41 (586 effort, 484 importance) and B86 (651 effort, 461 importance) would use 1237 effort (total 3096), also leaving 585 effort.} \\
\\ 
\texttt{Proposed counteroffer:} \\
\texttt{{-} Keep A87, C00, B32, C54 (effort = 4 + 116 + 285 + 209 = 614, importance = 668 + 121 + 583 + 508 = 1880).} \\
\texttt{{-} Replace A71 and A10 with B35, C01, B41 (effort = 664 + 571 + 586 = 182} \\
            }
        }
    }
     \\ \\

    \theutterance \stepcounter{utterance}  
    & & \multicolumn{4}{p{0.6\linewidth}}{
        \cellcolor[rgb]{0.9,0.9,0.9}{
            \makecell[{{p{\linewidth}}}]{
                \texttt{\tiny{[GM$\rangle$P2]}}
                \texttt{Your response did not start with the proper strategic reasoning tag at the very beginning of your response. The very first structured format must be of the form STRATEGIC REASONING: \{...\}. Try again.} \\
            }
        }
    }
    & \\ \\

    \theutterance \stepcounter{utterance}  
    & & & \multicolumn{4}{p{0.6\linewidth}}{
        \cellcolor[rgb]{0.9,0.9,0.9}{
            \makecell[{{p{\linewidth}}}]{
                \texttt{\tiny{[GM$\langle$P2]}}
                \texttt{STRATEGIC REASONING: \{} \\
\texttt{Opponent's proposal: \{'A87', 'C00', 'B32', 'C54', 'A71', 'A10'\} (effort = 1859, my importance = 2083).} \\
\texttt{Key insights:} \\
\texttt{1. Their proposal is under the LIMIT (3681) with 1822 effort remaining, allowing room for negotiation.} \\
\texttt{2. A87 (4 effort, 668 importance) is a critical inclusion due to its extreme efficiency.} \\
\texttt{3. A71 (505 effort, 65 importance) and A10 (740 effort, 138 importance) are low{-}value for me and should be replaced.} \\
\texttt{4. High{-}importance items missing: B35 (664 effort, 262 importance), C01 (571 effort, 389 importance), B41 (586 effort, 484 importance), B86 (651 effort, 461 importance), B33 (613 effort, 780 importance).} \\
\\ 
\texttt{Strategy:} \\
\texttt{{-} Propose a counteroffer that keeps A87, C00, B32, C54 (effort = 614, importance = 1880) and replaces A71/A10 with higher{-}value items.} \\
\texttt{{-} Target B35, C01, and B41 (total effort = 664 + 571 + 586 = 1821), which fits within remaining effort (1822) and adds 1135 importance (vs. 203 for A71/A10).} \\
\texttt{{-} This would yield total effort = 614 + 1821 = 2435 (<3681) and my importance = 1880 + 1135 = 3015.} \\
\texttt{\}} \\
\\ 
\texttt{PROPOSAL: \{'A87', 'C00', 'B32', 'C54', 'B35', 'C01', 'B41'\}} \\
\texttt{ARGUMENT: \{} \\
\texttt{Your proposal was a strong starting point, but I suggest replacing A71 and A10 with B35, C01, and B4} \\
            }
        }
    }
     \\ \\

    \theutterance \stepcounter{utterance}  
    & & & \multicolumn{2}{p{0.3\linewidth}}{
        \cellcolor[rgb]{0.9,0.9,0.9}{
            \makecell[{{p{\linewidth}}}]{
                \texttt{\tiny{[GM$|$GM]}}
                \texttt{abort} \\
            }
        }
    }
    & & \\ \\

    \theutterance \stepcounter{utterance}  
    & & & \multicolumn{2}{p{0.3\linewidth}}{
        \cellcolor[rgb]{0.9,0.9,0.9}{
            \makecell[{{p{\linewidth}}}]{
                \texttt{\tiny{[GM$|$GM]}}
                \texttt{end game} \\
            }
        }
    }
    & & \\ \\

\end{supertabular}
}

\end{document}
